\begin{afterword}
\summaryhead

The post argues that the laws of probability always govern reasoning, but that people usually cannot compute them and should not make up numerical probabilities from gut feelings. The brain has built-in ways of handling uncertainty, as in catching a ball, and a feeling put into words as ``one in a thousand'' will not be calibrated. Numbers should come from numbers.

The example is an argument given at the 2008 Global Catastrophic Risks conference: the safety papers on the Large Hadron Collider might be wrong, with a probability of at least 1 in 1000, and if they were, the world might be destroyed, again with at least 1 in 1000. The author objects to the air of authority of these numbers, proposes instead to debate a general rule of banning physics experiments that cannot be proved safe, and calls the speakers' call for more analysis disingenuous, since no paper could remove the doubt.

The author then admits an inconsistency: less alarm about the LHC than about a lottery device with a one-in-a-million chance of destroying the world, yet no confidence of being wrong only once in a million statements as sure as ``The Large Hadron Collider will not destroy the world''. The author keeps the inconsistency rather than make up numbers, since accuracy and winning matter more than consistency.

\responsehead

The general caution is sound and has support. Exact probabilistic inference is intractable in general, under the conditions the post states. The motor system combines prior knowledge with sensory uncertainty in a way close to Bayesian, without words. People who state extreme odds are overconfident: in a study the author's own chapter quotes, answers given odds of 1000 to 1 were right 81 per cent of the time. The post also limits its rule with care, to numbers produced by a non-numerical procedure, and it states its own inconsistency openly.

The example is handled with less care. The speaker's two figures were not equally baseless. The second, the chance of catastrophe if the safety argument fails, was a guess, as the speakers' written paper admits. The first, the chance that such an argument is flawed, was supported at the talk, by a journalist's report, with rates of retraction of journal articles. That is a number from numbers in the post's own sense, though the rates come from biomedicine. And the post's own test, put in the frequencies it recommends, places the author's chance of being wrong about the LHC above one in a million, while the speaker's figures multiply to at least one in a million. On the size of the risk the two are close.

The treatment of the speakers goes beyond the evidence. The post calls their call for more analysis ``disingenuous'' because no paper can remove the doubt; that reason would show at most a mistake, not bad faith. In their written paper, three months later, the speakers grant that the doubt cannot be removed and argue that independent arguments shrink it. The post's claim that a paper lowering the probabilities could only create ``the illusion'' of a different decision is asserted without argument. The general rule it proposes to debate is also broader than the one the speakers defended, as the author's own post of two days earlier indicates.

The closing preference for ``maybe'' and ``probably'' takes a side in an old debate. Sherman Kent of the CIA reported in 1964 that a board agreeing on ``serious possibility'' had meant odds from 20 to 80 to 80 to 20, and expected opponents to argue that numbers would make a writer look silly. The post argues from the writer's accuracy and leaves aside the reader.

In short: a sound caution against treating numbers made up from feelings as calibrated, applied to the LHC debate with less care: one of the speaker's two figures had a numerical basis, the author's own frequency test gives a risk of the same size, and the charge of disingenuousness is not supported.
\end{afterword}
